\documentclass{article}

\usepackage{amsmath,amssymb}
\usepackage{xurl}
\usepackage[hidelinks]{hyperref}

% Shared commands for notes in docs/paper-gaps/.

\DeclareUnicodeCharacter{2080}{\ifmmode{}_{0}\else\textsubscript{0}\fi}
\DeclareUnicodeCharacter{2081}{\ifmmode{}_{1}\else\textsubscript{1}\fi}
\DeclareUnicodeCharacter{2082}{\ifmmode{}_{2}\else\textsubscript{2}\fi}
\DeclareUnicodeCharacter{2099}{\ifmmode{}_{n}\else\textsubscript{n}\fi}

\newcommand{\Lean}{\textsc{Lean}}
\ExplSyntaxOn
\NewDocumentCommand{\leanid}{m}{
  \ifmmode
    \textsf{\detokenize{#1}}
  \else
    \group_begin:
    \urlstyle{sf}
    \tl_set:Nn \l_tmpa_tl {#1}
    \tl_replace_all:Nnn \l_tmpa_tl {\_} {_}
    \exp_args:NV \path \l_tmpa_tl
    \group_end:
  \fi
}
\ExplSyntaxOff
\providecommand{\tr}{\operatorname{tr}}
\newcommand{\tnleanIssueBase}{https://github.com/LionSR/TNLean/issues/}
\newcommand{\tnleanPrBase}{https://github.com/LionSR/TNLean/pull/}
\newcommand{\tnleanDocBase}{https://sirui-lu.com/TNLean/docs/}
\newcommand{\ghissue}[1]{\href{\tnleanIssueBase#1}{GitHub issue \##1}}
\newcommand{\ghpr}[1]{\href{\tnleanPrBase#1}{PR \##1}}
\newcommand{\doclink}[1]{\href{\tnleanDocBase#1.html}{\path{#1}}}

% Dirac notation and the identity operator, as in blueprint/src/macros/common.tex.
% \providecommand keeps note-local definitions authoritative.
\usepackage{dsfont}
\providecommand{\ket}[1]{|#1\rangle}
\providecommand{\bra}[1]{\langle #1|}
\providecommand{\braket}[2]{\langle #1 | #2 \rangle}
\providecommand{\ketbra}[2]{|#1\rangle\!\langle #2|}
\providecommand{\Id}{\mathds{1}}
\providecommand{\one}{\mathds{1}}
\providecommand{\C}{\mathbb{C}}
\providecommand{\MN}[1]{M_{#1}(\C)}
\providecommand{\MD}{M_D(\C)}

% External referencing for paper sources.  The build helper writes sanitized
% auxiliary files under build/paper-aux/ and excludes duplicated source labels.
\usepackage{xr}

% Import a generated paper auxiliary file with a caller-supplied label prefix.
% The candidate paths support builds from docs/paper-gaps/, the repository root,
% and build/paper-aux/ itself.
\newcommand{\importpaperaux}[2]{%
  \IfFileExists{../../build/paper-aux/#2.aux}{%
    \externaldocument[#1]{../../build/paper-aux/#2}%
  }{%
    \IfFileExists{build/paper-aux/#2.aux}{%
      \externaldocument[#1]{build/paper-aux/#2}%
    }{%
      \IfFileExists{#2.aux}{%
        \externaldocument[#1]{#2}%
      }{}%
    }%
  }%
}

% General external reference macros
\newcommand{\paperref}[2]{\ref{#1-#2}}
\newcommand{\papereqref}[2]{(\ref{#1-#2})}

% CPSV16 source (1606.00608 / MPDO-22-12-17-2.tex) external document and shortcuts
\importpaperaux{cpsv16-}{1606.00608/MPDO-22-12-17-2-tnlean}
\newcommand{\cpsvref}[1]{\paperref{cpsv16}{#1}}
\newcommand{\cpsveqref}[1]{\papereqref{cpsv16}{#1}}

% Verdict marker: \gapnote{<kind>}{<status>}. Kind is the policy
% classification (clarification, local-correction, scope-restriction,
% unfaithful, false-source, open-gap); status is open, wip (elimination
% underway), resolved, or historical. Machine-read by the site index;
% prints a verdict line.
\newcommand{\gapnote}[2]{\par\noindent\textbf{Verdict:} #1 (#2).\par}


\title{Reciprocal Factors in the Block-Independence Argument}
\author{}
\date{2026-08-29}

\begin{document}
\maketitle
\gapnote{local-correction}{resolved}

\section{The gauge multiplier}

Let $G$ act on a set $\mathsf X$, and let
$L\colon\mathsf X\times G^2\to\mathbb C^\times$ be an $L$-symbol. A fusion
scalar $\beta_{g,h}$ and an action scalar $\gamma^x_g$ act by multiplication
with
\begin{align}
    M^x_{g,h}(\beta,\gamma)
        &=\frac{\gamma^x_{gh}\beta_{g,h}}
        {\gamma^{h x}_g\gamma^x_h}.
    \label{eq:fbc25_bi_joint_gauge_multiplier}
\end{align}
At line~6990, Ref.~\cite{FrancoRubio2025MPUGauging} instead repeats
$\gamma_g$ in the second denominator factor. The second factor must be
$\gamma_h^x$, as in
(\ref{eq:fbc25_bi_joint_gauge_multiplier}); its group index is the second
argument of $L^x_{g,h}$.

\section{Factorization and the relative character}

The source defines
\begin{align}
    \ell^x_{a,b;r}
        &=\frac{L^x_{ar,r^{-1}b}}{L^x_{a,b}}.
    \label{eq:fbc25_bi_ell_definition}
\end{align}
Taking $a=g$, $b=h$, and $r=h$ in
(\ref{eq:fbc25_bi_ell_definition}) gives
\begin{align}
    \ell^x_{g,h;h}
        &=\frac{L^x_{gh,e}}{L^x_{g,h}},
    \label{eq:fbc25_bi_ell_specialization}\\
    L^x_{g,h}
        &=\frac{L^x_{gh,e}}{\ell^x_{g,h;h}}.
    \label{eq:fbc25_bi_ell_solved}
\end{align}
Thus line~6996 of Ref.~\cite{FrancoRubio2025MPUGauging} must divide by
$\ell^x_{g,h;h}$ rather than multiply by it. Compatibility at
$(gh,e,e)$ gives
\begin{align}
    L^x_{gh,e}
        &=\frac{L^x_{e,e}}{\omega(gh,e,e)},
    \label{eq:fbc25_bi_right_identity}\\
    L^x_{g,h}
        &=\frac{L^x_{e,e}}
        {\omega(gh,e,e)\ell^x_{g,h;h}}.
    \label{eq:fbc25_bi_corrected_factorization}
\end{align}
Under block independence, $\ell^x_{g,h;h}$ may be evaluated at any fixed
$x_0\in\mathsf X$. Hence the corrected block factor is
\begin{align}
    A(g,h)
        &=(\omega(gh,e,e)\ell^{x_0}_{g,h;h})^{-1},
    \label{eq:fbc25_bi_block_factor}\\
    L^x_{g,h}
        &=A(g,h)L^x_{e,e}.
    \label{eq:fbc25_bi_factorized_lsymbol}
\end{align}
These identities are formalized directly from the definition of $\ell$ and
compatibility.\footnote{Formal declarations:
\leanid{TNLean.Algebra.LSymbol.ell},
\leanid{TNLean.Algebra.LSymbol.blockFactor}, and
\leanid{TNLean.Algebra.LSymbol.factorization_of_isBlockIndependentEll} in
\doclink{TNLean/Algebra/LSymbolBlockIndependence}.}

Compatibility at $(e,e,g)$ also gives
\begin{align}
    L^x_{e,g}
        &=\omega(e,e,g)L^{g x}_{e,e}.
    \label{eq:fbc25_bi_left_identity}
\end{align}
Combining (\ref{eq:fbc25_bi_corrected_factorization}) at $(e,g)$ with
(\ref{eq:fbc25_bi_left_identity}) yields
\begin{align}
    L^{g x}_{e,e}
        &=\rho(g)L^x_{e,e},
    \label{eq:fbc25_bi_relative_transport}\\
    \rho(g)
        &=\bigl(\ell^{x_0}_{e,g;g}\omega(g,e,e)
            \omega(e,e,g)\bigr)^{-1}.
    \label{eq:fbc25_bi_corrected_relative_character}
\end{align}
The factor $\ell^{x_0}_{e,g;g}$ at line~7016 of
Ref.~\cite{FrancoRubio2025MPUGauging} must therefore occur reciprocally.
Applying (\ref{eq:fbc25_bi_relative_transport}) successively to $h$ and $g$
shows that
\begin{align}
    \rho(gh)
        &=\rho(g)\rho(h).
    \label{eq:fbc25_bi_character_multiplicativity}
\end{align}
No transitivity or finiteness assumption on the action is needed.\footnote{Formal
declarations: \leanid{TNLean.Algebra.LSymbol.relativeCharacter},
\leanid{TNLean.Algebra.LSymbol.relativeCharacter_mul_base}, and
\leanid{TNLean.Algebra.LSymbol.relativeCharacter_mul} in
\doclink{TNLean/Algebra/LSymbolBlockIndependence}.}

\section{Forward and inverse gauges}

Define
\begin{align}
    \beta^+_{g,h}
        &=\frac{A(g,h)}{\rho(h)},
    \label{eq:fbc25_bi_forward_fusion_gauge}\\
    \gamma^{+,x}_g
        &=(L^{g x}_{e,e})^{-1}.
    \label{eq:fbc25_bi_forward_action_gauge}
\end{align}
Using (\ref{eq:fbc25_bi_relative_transport}), their multiplier is
\begin{align}
    M^x_{g,h}(\beta^+,\gamma^+)
        &=\frac{A(g,h)}{\rho(h)}
          \frac{(L^{ghx}_{e,e})^{-1}}
          {(L^{ghx}_{e,e})^{-1}(L^{hx}_{e,e})^{-1}}
    \label{eq:fbc25_bi_forward_multiplier_first}\\
        &=A(g,h)L^x_{e,e}
    \label{eq:fbc25_bi_forward_multiplier_second}\\
        &=L^x_{g,h}.
    \label{eq:fbc25_bi_forward_multiplier_lsymbol}
\end{align}
Consequently, $(\beta^+,\gamma^+)$ sends the trivial $L$-symbol to $L$:
\begin{align}
    (\beta^+,\gamma^+)\mathbin{\triangleright}1
        &=L.
    \label{eq:fbc25_bi_forward_gauge_produces_lsymbol}
\end{align}
The displayed factor in the source is therefore a forward gauge multiplier,
not a multiplier that trivializes $L$. The pointwise inverse gauges give the
required conclusion:
\begin{align}
    ((\beta^+)^{-1},(\gamma^+)^{-1})
        \mathbin{\triangleright}L
        &=1.
    \label{eq:fbc25_bi_inverse_gauge_trivializes_lsymbol}
\end{align}
For a supplied compatible scalar $L$-symbol, this proves the implication from
block independence to gauge triviality and, together with the immediate
converse implications, the four-way scalar equivalence used in the proof of
Proposition~\texttt{prop:technical01} of
Ref.~\cite{FrancoRubio2025MPUGauging}.\footnote{Formal declarations:
\leanid{TNLean.Algebra.LSymbol.factorFusionGauge},
\leanid{TNLean.Algebra.LSymbol.factorActionGauge},
\leanid{TNLean.Algebra.LSymbol.gauge_one_eq_of_isBlockIndependentEll},
\leanid{TNLean.Algebra.LSymbol.gauge_inv_eq_one_of_isBlockIndependentEll}, and
\leanid{TNLean.Algebra.LSymbol.hasTrivialGauge_of_isBlockIndependentEll} in
\doclink{TNLean/Algebra/LSymbolBlockIndependence}. The capstone equivalences
are \leanid{TNLean.Algebra.LSymbol.hasTrivialGauge_iff_hasBlockConstantGauge},
\leanid{TNLean.Algebra.LSymbol.hasTrivialGauge_iff_hasTrivialEllGauge}, and
\leanid{TNLean.Algebra.LSymbol.hasTrivialGauge_iff_isBlockIndependent}.}

\section{Scope of the formal result}

The formal declarations start with a supplied compatible scalar $L$-symbol and
prove the algebraic equivalence among its four scalar gauge properties. The
fourth condition of Proposition~\texttt{prop:technical01} in
Ref.~\cite{FrancoRubio2025MPUGauging} is instead stated for an MPS--MPU pair.
A tensor-facing result must still construct the associated $L$-symbol and prove
that the MPS and MPU tensor gauges induce the scalar multiplier
(\ref{eq:fbc25_bi_joint_gauge_multiplier}). The scalar calculation does not
supply that bridge and therefore does not by itself establish the full source
proposition.

\section{Conclusion}

The four local corrections form one reciprocal chain. The denominator of the
joint gauge multiplier uses $\gamma_h^x$; the definition of $\ell$ forces
division by $\ell^x_{g,h;h}$ in the factorization; the relative character
consequently contains the reciprocal of $\ell^{x_0}_{e,g;g}$; and the resulting
factor is a forward gauge multiplier whose inverse trivializes the supplied
$L$-symbol. These reciprocal corrections are resolved. The tensor-facing bridge
required for the full statement of Proposition~\texttt{prop:technical01}
remains outside the scalar result.

\bibliographystyle{alpha}
\bibliography{references}

\end{document}